\section{Related Work} \label{sec:related_work}
Approaches in continual learning can be grouped into three main categories. We review them below.


\textbf{Regularization-based}:
This approach computes importance values of either parameters or their gradients on previous tasks, and adds a regularization in the loss to restrict changes to those important parameters to mitigate CF. 
EWC~\citep{ewc} uses the Fisher information matrix to represent the importance of parameters and a regularization to penalize the sum of changes to all parameters. 
SI~\citep{si} extends EWC to reduce the complexity in computing the penalty.
Many other approaches~\citep{lwf,dmc,ucl} in this category have also been proposed, but they still have difficulty to prevent CF.  
As discussed in the introduction section, the proposed approach SPG has some resemblance to a  regularization based method EWC. But the coarse-grained approach of using regularization is significant poorer than the fine-grained soft-masking in SPG for overcoming CF as we will see in \cref{sec:results}.



\textbf{Memory-based}: 
This approach introduces a small memory buffer to store data of previous tasks and replay them in learning a new task to prevent CF~\citep{gem,a-gem}. Some methods~\citep{gr,binplay} prepare data generators for previous tasks, and the generated pseudo-samples are used instead of real samples. 
Although several other approaches~\citep{icarl,mer,mir} have been proposed, they still suffer from CF. SPG does not save replay data or generate pseudo-replay data.



\textbf{Parameter isolation-based}:
This approach is most similar to ours SPG.
It tries to learn a sub-network for each task (tasks may share parameters and neurons), which limits knowledge transfer. We have discussed HAT, SupSup, and CAT in \cref{sec:introduction}.
Many others also take similar approaches, e.g., 
Progressive Networks (PGN)~\citep{pgn}, APD~\citep{apd}, PathNet~\citep{pathnet}, PackNet~\citep{packnet}, SpaceNet~\citep{spacenet}, and WSN~\citep{wsn}. 
In particular, PGN allocates a sub-network for each task in advance, and progressively concatenates previous sub-networks while freezing parameters allocated to previous tasks.
APD selectively reuses and dynamically expands the dense network.
Those methods, however, depend on the expansion of network for their performance, which is often not acceptable in cases where many tasks need to be learned.
PathNet splits each layer into multiple sub-modules and finds the best \textit{pathway} designated to each task. 
PackNet freezes important weights for each task by finding them based on pruning.
Although PathNet and PackNet do not expand 
the network along with continual learning, they suffer from  over-consumption of the fixed capacity.
To address this, SpaceNet adopts the sparse training to preserve parameters for future tasks but the performance for each task is sacrificed.
WSN also allocates a subnetwork within a dense network
and selectively reuses subnetworks for previous tasks without expanding the whole network.
Nevertheless, those methods are still limited by the pre-allocated network size as each task monopolizes and consumes some amount of capacity, which results in poorer KT when learning many tasks.


In summary, parameter isolation-based methods suffer from over-consumption of network capacity and have limited KT, which the proposed method tries to address at the same time.



